\scanpage{8}
\mspage{4}
\uline{NB}
\[
\inserted{D(M)=}
\operatorname{Hom}_A(M,U_g)
=\operatorname{Hom}_A(M_g,\operatorname{Hom}_{k'}(A_g,k'))
\]
\[
\simeq\deleted{\unsure{A_g}}\quad
\text{module des homs }k\text{-bil.}\quad
A_g\times M_g\overset{\varphi}{-}k'
\]
telles que
\[
\begin{aligned}
\varphi(\lambda u,x)&=\lambda\varphi(u,x)\\
\varphi(u\deleted{\unsure{\otimes}}v,x)&=\varphi(u,v.x)
\end{aligned}
\qquad
\begin{gathered}
\lambda\in k'\\u\in A\\x\in\unsure{u|M}
\end{gathered}
\]
i.e. correspondent aux homs
\[
A_d\otimes_A M\longrightarrow k'
\]
qui sont $k'$-linéaires, pour la structure de $k'$-module à gauche
provenant de $A_d$.

Or
\[
\begin{aligned}
A_d\otimes_A M&\xleftarrow{\sim}M
&&\text{isom de }k'\text{-modules}\\
1\otimes x&\mathrel{\longleftarrow\!\vert} x\\
\lambda(1\otimes x)=\lambda\otimes x=1\otimes\lambda x
&\mathrel{\longleftarrow\!\vert}\lambda x
\end{aligned}
\]
donc on trouve bien \deleted{\unsure{$M^{\vee(k')}$}}
\[
F'(D(M))=M^{\vee(k')}
\]
et la donnée de $a:A-A^o$ permet de mettre sur
$M^{\vee(\unsure{k'|k})}$ une structure de $A$-module.

Il faut encore que
\[
\deleted{\unsure{M\otimes M^{\vee(k')}}}
\]
\[
\underbrace{F'(M)}_{\substack{\deleted{\unsure{F'(M)}}\\
M_{\deleted{\unsure{k'}}}^{k'}}}
\otimes
\underbrace{F'(D(M))}_{(M^{k'})^\vee}
\longrightarrow k'=F'(1_C)
\]
provient bien \deleted{$A$} d'un hom
\[
M\otimes D(M)\longrightarrow1_C
\]
\margin{7}
